%\SweaveUTF8




\setkeys{Gin}{width = \textwidth}

\begin{Problem}[Sample size]\label{problem:sample_size}
We use the data set \code{sepe.dat} from Problem Set 1.

\begin{Subtask}
    %---------------------------------------------------------------------------
    \item We plan a new sample survey to estimate total environmental
        protection expenditures. The survey will use simple random sampling
        (\SRS), and the total will be estimated by $\tys = N\mys$, where $\mys$
        is the sample mean.

        For sample size planning, assume that future expenditures have
        approximately the same coefficient of variation as the expenditures
        observed in \code{sepe.dat}. Determine the sample size so that the
        relative error exceeds $\rho = 10\%$ with probability at most $\alpha =
        5\%$.

        For the population mean, the requirement is
        \begin{equation*}
            \PP\left(\left\vert \frac{\mys - \myu}{\myu} \right\vert >
                \rho\right) \leq \alpha.
        \end{equation*}
        Rewrite this condition for $\tys$. What is the minimum required sample
        size?

    \begin{Solution}
        Since $\tys = N\mys$ and $\ty = N \myu$,
        \begin{equation*}
            \left\vert \frac{\tys - \ty}{\ty} \right\vert =
            \left\vert \frac{\mys - \myu}{\myu} \right\vert.
        \end{equation*}
        Thus, the same relative-error condition applies to the estimator of the
        total. Using a normal approximation and initially ignoring the
        finite-population correction, we obtain
        \begin{equation*}
            n_0 = \left(\frac{z_{1 - \alpha/2} \cdot
            \sigmay}{\rho\,\myu}\right)^2 = \left(\frac{z_{1 -
            \alpha/2}\cdot\mathrm{CV}_y}{\rho}\right)^2,
        \end{equation*}
        where $\sigmay$ is the finite-population standard deviation and
        $\mathrm{CV}_ = \sigmay/\myu$ is the finite-population coefficient of
        variation. Accounting for the finite-population correction gives
        \begin{equation*}
            n = \frac{n_0}{1 + n_0/N}.
        \end{equation*}

        \textbf{Calculation in R:}


\begin{Schunk}
\begin{Sinput}
> sepe <- read.table("https://stat.ethz.ch/Teaching/Datasets/cas-das/sepe.dat",
                     header = TRUE)
> attach(sepe)
\end{Sinput}
\end{Schunk}

\begin{Schunk}
\begin{Sinput}
> alpha <- 0.05
> rho <- 0.1
> N <- NROW(sepe)
> CV_y <- sd(expend) / mean(expend)
> z <- qnorm(1 - alpha / 2)
\end{Sinput}
\end{Schunk}

        Ignoring the finite population correction,

\begin{Schunk}
\begin{Sinput}
> n0 <- (z * CV_y / rho)^2
> ceiling(n0)
\end{Sinput}
\begin{Soutput}
[1] 242
\end{Soutput}
\end{Schunk}

        Including the finite population correction,

\begin{Schunk}
\begin{Sinput}
> n <- n0 / (1 + n0 / N)
> ceiling(n)
\end{Sinput}
\begin{Soutput}
[1] 186
\end{Soutput}
\end{Schunk}

        Hence, the sample should contain at least 186 units.
    \end{Solution}

    %---------------------------------------------------------------------------
    \item Use the minimum sample size from part (a) to estimate total
        environmental protection expenditures and report a 95\% confidence
        interval. Cochran proposed the rule $n > 25 G_1^2$, where $G_1$ is
        Fisher's coefficient of skewness, i.e., the third standardized
        moment.\footnote{Cochran, W. G. (1977). Sampling Techniques.
            3$^{\mathrm{rd}}$ ed., New York: John Wiley \& Sons, p. 42.}
        Sugden et al. proposed the modified rule $n > 28 + 25
        G_1^2$.\footnote{Sugden, R. A., Smith, T. M. F., and Jones, R. P.
            (2000). Cochran's Rule for Simple Random Sampling. Journal of the
            Royal Statistical Society Series B, 62, pp. 787--793.}

        Does the sample size satisfy Cochran's rule? Does it satisfy the rule
        proposed by Sugden et al.? What do these rules suggest about the
        adequacy of the 95\% confidence interval?

        \begin{Rhint}
            You can use the following function to compute the finite-population
            moment coefficient of skewness.

\begin{Scode}
skewness_pop <- function(x)
{
    z <- x - mean(x)
    mean(z^3) / mean(z^2)^(3 / 2)
}
\end{Scode}

        \end{Rhint}

    \begin{Solution}
        The skewness of \code{expend} is approximately


\begin{Schunk}
\begin{Sinput}
> G1 <- skewness_pop(expend)
> G1
\end{Sinput}
\begin{Soutput}
[1] 2.276473
\end{Soutput}
\end{Schunk}

        Cochran's rule requires
        \begin{equation*}
            n > 25 G_1^2,
        \end{equation*}
        which gives

\begin{Schunk}
\begin{Sinput}
> ceiling(25 * G1^2)
\end{Sinput}
\begin{Soutput}
[1] 130
\end{Soutput}
\end{Schunk}

        Sugden et al.'s rule for a Studentized confidence interval is
        \begin{equation*}
            n > 28 + 25 G_1^2,
        \end{equation*}
        giving

\begin{Schunk}
\begin{Sinput}
> ceiling(28 + 25 * G1^2)
\end{Sinput}
\begin{Soutput}
[1] 158
\end{Soutput}
\end{Schunk}

        The planned sample size of 186 satisfies both rules.
        Thus, these rules support using the usual normal-based 95\,\% confidence
        interval. They do not, however, guarantee an exact 95\,\% coverage
        probability.
    \end{Solution}

    %---------------------------------------------------------------------------
    \item Under the same assumptions as in part a), determine the sample size so
        that the coefficient of variation of the estimator
        \begin{equation*}
            \varcoef(\,\tys) = \frac{\se(\,\tys)}{\EE(\,\tys)}
        \end{equation*}
        is at most 5\,\%, where $\tys = N \mys$.

    \begin{Solution}
        Under \SRS,
        \begin{equation*}
            \var(\,\tys) = N^2\left(1-\frac{n}{N}\right)\frac{\varyu}{n},
            \qquad \EE(\,\tys)=\ty=N\myu.
        \end{equation*}
        Hence,
        \begin{equation*}
            \varcoef(\,\tys) = \mathrm{CV}_y \sqrt{\frac{1-n/N}{n}}.
        \end{equation*}

        We require
        \begin{equation*}
            \varcoef(\,\tys) \leq 0.05.
        \end{equation*}
        Ignoring the finite-population correction first gives
        \begin{equation*}
            \frac{\mathrm{CV}_y}{\sqrt{n}} \leq 0.05,
        \end{equation*}
        or equivalently,
        \begin{equation*}
            n \geq n_0 = \left(\frac{\mathrm{CV}_y}{0.05}\right)^2.
        \end{equation*}

        Accounting for the finite-population correction gives
        \begin{equation*}
            n \geq \frac{n_0}{1+n_0/N}.
        \end{equation*}
        Thus, the minimum sample size is the smallest integer satisfying this
        inequality.

        \textbf{Calculation in R:}

\begin{Schunk}
\begin{Sinput}
> n0 <- (CV_y / 0.05)^2
> ceiling(n0)
\end{Sinput}
\begin{Soutput}
[1] 252
\end{Soutput}
\end{Schunk}

\begin{Schunk}
\begin{Sinput}
> n <- n0 / (1 + n0 / N)
> ceiling(n)
\end{Sinput}
\begin{Soutput}
[1] 192
\end{Soutput}
\end{Schunk}

        Thus, at least 192 units are required.
    \end{Solution}

    %---------------------------------------------------------------------------
    \item How large must the sample be to estimate the population proportion
        $\propu_{31}$ of establishments in industry class 31 (indicator
        $W_{31}$) with a standard error of at most 5 percentage points? Use
        \begin{equation*}
            \props_{31} = \frac{1}{n}\sum_{i \in s}
                \mathbbm{1}\{i \in W_{31}\},
        \end{equation*}
        where
        \begin{equation*}
            \indicator\{i \in W_{31}\} = \begin{cases}
                1 & \text{if $i$ belongs to industry class 31},\\
                0 & \text{otherwise}.
            \end{cases}
        \end{equation*}

    \begin{Solution}
        From Problem 1.2c), $\props_{31} = 0.1725$. Let $\indicator_{31}$
        denote the indicator variable defined above. Under \SRS,
        \begin{equation*}
            \se(\,\props_{31}) =
            \sqrt{\left(1 - \frac{n}{N}\right)\frac{\varu_{\indicator_{31}}}{n}},
        \end{equation*}
        where $\varu_{\indicator_{31}}$ is the finite-population variance of
        the indicator.

        Ignoring the finite population correction first,
        \begin{equation*}
            n_0 = \frac{\varu_{\indicator_{31}}}{0.05^2}.
        \end{equation*}
        The corrected sample size is
        \begin{equation*}
            n = \frac{n_0}{1 + n_0/N}.
        \end{equation*}

        \textbf{Calculation in R:}

\begin{Schunk}
\begin{Sinput}
> I31 <- code
> P31 <- mean(I31)
> sigma2_I31 <- var(I31)
> n0 <- sigma2_I31 / 0.05^2
> ceiling(n0)
\end{Sinput}
\begin{Soutput}
[1] 58
\end{Soutput}
\end{Schunk}

\begin{Schunk}
\begin{Sinput}
> n <- n0 / (1 + n0 / N)
> ceiling(n)
\end{Sinput}
\begin{Soutput}
[1] 54
\end{Soutput}
\end{Schunk}

        Thus, the sample should contain at least 54
        establishments.
    \end{Solution}
\end{Subtask}
\end{Problem}
